\chapter{NeurIPS 实验设计与评价方法}
\label{chap:experiment-design}

本章把当前实验协议升级为 NeurIPS main-track 的监督实验矩阵。核心原则是：每一个实验只服务一个可证伪命题，每一个主表数值都必须由 logger、原始输出、统计脚本和 checksum 共同支撑。现有 Protocol-A/Protocol-C 仍是基线证据，但它们只能证明“当前已评估检索基线下 TRR 有优势”；若要达到主会方法论文强度，还必须补齐强基线、消融、泄漏审计、真实鲁棒性和感知相关性。

\begin{table}[H]
	\centering
	\caption{NeurIPS main-track 补实验矩阵与 KaiMing 通过标准。}
	\label{tab:neurips-experiment-matrix}
	\scriptsize
	\resizebox{\textwidth}{!}{%
	\begin{tabular}{p{0.08\textwidth}p{0.19\textwidth}p{0.25\textwidth}p{0.22\textwidth}p{0.18\textwidth}}
		\toprule
		实验 & 目标命题 & 最小方法集合 & 主要指标 & 通过标准 \\
		\midrule
		E0 Evidence freeze & 当前主证据可复核 & Protocol-A 204、Protocol-C、listening test & checksum、run manifest、统计复现 & 所有主表文件可由 manifest 定位，checksum 通过。 \\
		E1 Strong baselines & 排除弱 baseline 解释 & Text, Wav2Vec, FeatureNN, CLAP, PaSST, PANNs, direct regression, Text2FX-style optimization & Param. Dist., Acc@0.1, Recall, Cosine, Module & 同 split、同 query set、同 logger 下报告 CI 与 Holm-corrected tests。 \\
		E2 TRR ablation & 二阶纹理统计是有效成分 & mean pooling, Gram statistics, layer choice, projection dim, normalization, top-k aggregation & 主指标加运行成本 & 至少一个结构性消融显著削弱主指标，说明贡献不只是实现偶然性。 \\
		E3 Leakage and duplicate audit & 结果不来自近重复泄漏 & exact path, perceptual hash, embedding nearest-neighbor, family/group split & duplicate rate, performance under stricter split & 更严格 split 下结论方向保持，或明确降级 claim。 \\
		E4 Real robustness & 排除 embedding-noise-only 解释 & clean, SNR noise, EQ, compression, reverberation, codec & 参数指标、fallback weight、失败类型 & fusion/TRR 边界可解释；不能只报告有利场景。 \\
		E5 Metric perception link & 参数距离与听感有关系 & objective metric vs listening score, rater metadata, anchor design & Spearman/Pearson, CI, inter-rater reliability & 若相关性弱，主文必须把指标定位为 controllability proxy。 \\
		E6 Submission package & 可复现与匿名合规 & code zip, data card, checklist, appendix, run scripts & checklist completeness, anonymous artifact audit & 论文 claim、附录表、代码输出三者一一对应。 \\
		\bottomrule
	\end{tabular}}
\end{table}

表~\ref{tab:neurips-experiment-matrix} 是后续 4--6 周的实验契约。该契约把当前最可辩护的研究对象压实为四个限定词：检索式、可执行、参数可编辑、证据可审计。凡是不能进入这个矩阵的探索，例如产品侧 UI 或端到端生成体验，都应保留在 future work 或系统背景中。

本报告采用三类协议口径。第一类是主证据协议，即 resolved-audio-grouped Protocol-A。该协议基于外部 1267 条 guitar-effects benchmark drop-in，按 resolved local audio path 分组后形成 204 条测试查询和 1063 条知识库记录，目的在于消除 legacy split 中已发现的 exact cross-split audio reuse。第二类是诊断协议，包括 211-query Protocol-C 和 legacy/strong-baseline 结果，用于分析融合边界和历史数值差异。第三类是待补协议，即 P0 E2/E3；它有报告和脚本，但缺少原始 outputs，因此不能作为主证据。

\begin{table}[H]
	\centering
	\caption{实验协议、数据规模与证据状态。}
	\label{tab:protocol-map}
	\small
	\resizebox{\textwidth}{!}{%
	\begin{tabular}{lrrll}
		\toprule
			协议或材料 & 测试数 & 知识库数 & 状态 & 本文用途 \\
		\midrule
		Protocol-A audio-grouped & 204 & 1063 & 本地主统计与 per-query 结果存在 & 主证据，论文主表口径。 \\
		Protocol-A legacy/strong & 211 & 1056 & 本地结果存在，但含 legacy split 风险 & 历史对照或诊断，不与 204 口径混写。 \\
		Protocol-C & 211 & 1056 & 本地统计与 per-query 结果存在 & Fusion fallback 与 conflict failure 诊断。 \\
		Listening test & 26 participants & 910 ratings & 统计报告存在 & 主观补充，不作为严格胜负证据。 \\
		P0 E2/E3 & 204 & 1063 & 报告和脚本存在，原始 outputs 缺失 & 复现入口，不能作为正式主证据。 \\
		Protocol-B & 211 & 1056 & 文件自标 deprecated & 历史材料，不进入主结论。 \\
		\bottomrule
	\end{tabular}}
\end{table}

表~\ref{tab:protocol-map} 中最重要的设计选择是把 204 条 audio-grouped Protocol-A 固定为主文口径。旧 211-query 结果在数值上更好看，例如 TRR 的 L2 为 0.3064，但该切分的 audit 发现 16 条跨 split shared resolved audio paths 和 48 对近重复风险。严格切分的 TRR L2 变为 8.0467，数值尺度明显不同，却更适合作为提交面主证据。因此，本文宁可使用较难、更保守的主表，也不将 legacy 数值当作核心结论。

每个检索方法输出一个候选参数 JSON。参数距离按展开后的数值叶子计算 RMSE；Acc@0.1、Recall、Cosine 和 Module 由 \texttt{Experiments/common/evaluate.py} 中的 \texttt{Evaluator} 实现。Protocol-A 统计报告使用 bootstrap confidence intervals 和 paired permutation tests，并用 Holm correction 处理多重比较。CLAP 在 audio-grouped 协议中只覆盖 201 条查询，因此与 TRR 的配对比较使用交集样本。

Protocol-C 的四个场景分别测试标准输入、vague text、noisy audio 和 modality conflict。需要强调的是，noisy audio 场景使用 embedding-space Gaussian noise，而非真实音频退化重新编码。因此，Protocol-C 只能说明当前 score-level fusion heuristic 的边界行为，不能证明系统已具备真实环境下的音频鲁棒性。
